\documentclass[11pt]{article}
\usepackage[margin=1in]{geometry}
\usepackage{graphicx}
\usepackage{booktabs}
\usepackage{float}
\usepackage{natbib}
\title{Topology-State Analysis of Random Peptide Interaction Networks}
\author{}
\date{}
\begin{document}
\maketitle
\begin{abstract}
We analyzed 1000 peptide sequences with paired monomer and dimer network representations and quantified graphlet-defined amyloid topology classes at the sequence level. Dominant topology distributions shifted between states, with `none` changing from 43.7\% (monomer) to 48.3\% (dimer) and `double 1,2 2-ribbon` from 31\% to 44.4\%. Transition analysis identified stable `none` and stable `double 1,2 2-ribbon` as the most frequent outcomes.
\end{abstract}
\section{Methods}
Edgelists were loaded from `random_sequences_edgelist.rda` (paired monomer/dimer networks). For each network, node-level topology membership was computed with the ORCA-based fibril assay used in prior Grazioli-Butts amyloid network work. Sequence-level topology fractions were computed as class-wise mean node membership, and the dominant class per sequence was assigned as the maximum-fraction class (or `none` if all class fractions were zero; `mixed` for ties).
\section{Results}
Each state included 1000 sequences. The most common dominant topology in monomer was `none` (43.7\%) followed by `double 1,2 2-ribbon` (31\%). In dimer, `none` (48.3\%) and `double 1,2 2-ribbon` (44.4\%) remained dominant.
The largest transition counts were `none \rightarrow none` (293, 29.3\%), `double 1,2 2-ribbon \rightarrow double 1,2 2-ribbon` (163, 16.3\%), `double 1,2 2-ribbon \rightarrow none` (137, 13.7\%), and `none \rightarrow double 1,2 2-ribbon` (131, 13.1\%).
\begin{figure}[H]
\centering
\includegraphics[width=0.95\linewidth]{../plots/topology_fraction_boxplots.png}
\caption{Distribution of sequence-level topology fractions by state.}
\end{figure}
\begin{figure}[H]
\centering
\includegraphics[width=0.85\linewidth]{../plots/dominant_topology_counts.png}
\caption{Dominant topology counts in monomer and dimer states.}
\end{figure}
\begin{figure}[H]
\centering
\includegraphics[width=0.78\linewidth]{../plots/topology_transition_heatmap.png}
\caption{Monomer-to-dimer transition matrix of dominant topology classes.}
\end{figure}
\begin{table}[ht]
\centering
\caption{Dominant topology composition in monomer and dimer ensembles (N=1000 sequences per state).}
\label{tab:dominant-composition}
\begin{tabular}{llrr}
\toprule
State & Dominant topology & Count & Percent (\%) \\
\midrule
dimer & none & 483 & 48.3 \\
dimer & double 1,2 2-ribbon & 444 & 44.4 \\
dimer & 1,2 2-ribbon & 43 & 4.3 \\
dimer & mixed & 27 & 2.7 \\
dimer & 1-ribbon & 3 & 0.3 \\
monomer & none & 437 & 43.7 \\
monomer & double 1,2 2-ribbon & 310 & 31.0 \\
monomer & 1,2 2-ribbon & 170 & 17.0 \\
monomer & mixed & 57 & 5.7 \\
monomer & 1-ribbon & 26 & 2.6 \\
\bottomrule
\end{tabular}
\end{table}

\begin{table}[ht]
\centering
\caption{Top monomer-to-dimer dominant-topology transitions.}
\label{tab:top-transitions}
\begin{tabular}{llrr}
\toprule
Monomer dominant & Dimer dominant & Count & Percent (\%) \\
\midrule
none & none & 293 & 29.3 \\
double 1,2 2-ribbon & double 1,2 2-ribbon & 163 & 16.3 \\
double 1,2 2-ribbon & none & 137 & 13.7 \\
none & double 1,2 2-ribbon & 131 & 13.1 \\
1,2 2-ribbon & double 1,2 2-ribbon & 106 & 10.6 \\
1,2 2-ribbon & none & 37 & 3.7 \\
mixed & double 1,2 2-ribbon & 37 & 3.7 \\
1,2 2-ribbon & 1,2 2-ribbon & 16 & 1.6 \\
mixed & none & 13 & 1.3 \\
1-ribbon & 1,2 2-ribbon & 12 & 1.2 \\
\bottomrule
\end{tabular}
\end{table}

\section{Additional States Metrics}
We computed network-level organization metrics on the same monomer/dimer edgelists as complementary descriptors of self-assembly patterns. These metrics do not define explicit topology classes (e.g., 1-ribbon/2-ribbon); instead, they summarize connectivity structure and its changes between states. Assortativity captures degree-similarity in contacts (more assortative implies clustered, like-degree mixing). The number of communities summarizes interaction compartmentalization. Core statistics (core max-$k$ and core fraction) summarize the strength and size of the mutual-connectivity scaffold. Role entropy and role polarization summarize heterogeneity in graphlet-orbit roles. Operational state flags binarize these metrics to report how frequently networks are assortative, modular, core-periphery-like, or role-polarized.
Across paired sequences, dimerization is associated with a small decrease in assortativity (median $\Delta=-0.018$), a rise in community count (median $\Delta=+1$), a marked reduction in core fraction (median $\Delta=-0.120$; core max-$k$ roughly stable), and a strong shift toward higher role entropy (median $\Delta=+0.307$) with lower role polarization (median $\Delta=-0.071$). Operational flags are largely unchanged for assortative/modular categories, while core-periphery flags increase modestly in the dimer state (mean $\Delta\approx+0.17$).
\begin{figure}[H]
\centering
\includegraphics[width=0.98\linewidth]{../plots/additional_states/additional_states_metric_distributions_ACS_doublecol.png}
\caption{Additional-states metrics with panel labels (a--n). Left column: raw monomer vs dimer distributions. Right column: per-sequence changes (dimer minus monomer). Panels (a,b) assortativity; (c,d) number of communities; (e,f) core max-$k$; (g,h) core fraction; (i,j) role entropy; (k,l) role polarization; (m,n) core--periphery flag.}
\end{figure}
\section{Reproducibility Files}
Core outputs used by this manuscript: `topology_fractions.csv`, `dominant_topology_counts.csv`, `transition_counts.csv`, and the PNG plots in `../plots`.
\bibliographystyle{plainnat}
\bibliography{references}
\end{document}
